\section{Program Design and Implementation}

\begin{figure}[H]
  \centering
  \includegraphics[height=0.72\textheight,keepaspectratio]{figures/stereo_pipeline.png}
  \caption{Stereo depth estimation pipeline generated from the Mermaid source file.}
  \label{fig:stereo-pipeline}
\end{figure}

The program entry point is \codepath{task2/src/stereo_depth.py}, which implements a single
reproducible run from input validation to result serialization. The program does not use display
images as numerical inputs; instead, it saves both the original floating-point arrays and the
8-bit visualizations used in the report. The overall processing relationship is shown in
Fig.~\ref{fig:stereo-pipeline}.

\subsection{Input Reading, Validation, and Preprocessing}

The program uses \texttt{cv2.imread} to read the left and right color images. Before matching,
it checks that the files exist, that OpenCV returns non-empty arrays, and that the two images have
the same row and column dimensions. If either image is empty or the dimensions are inconsistent,
the program raises a locatable error and does not continue to produce seemingly valid results.

The two images are then converted to grayscale. To reduce the computational cost, the left and
right color images and their grayscale versions are downsampled once with
\texttt{cv2.pyrDown}. Thus, no different scaling ratio is introduced between the two views. The
original size of $1282\times1110$ becomes $641\times555$; all subsequent disparity arrays and
output images correspond to this processed size.

\subsection{StereoSGBM Configuration}

The parameters used in the actual run are recorded in \codepath{task2/results/params.json} and
listed in Table~\ref{tab:sgbm-params}. \texttt{numDisparities} is set to a multiple of 16 so
that the search range covers the main disparity variation in the official aloe images.

\begin{table}[H]
  \centering
  \caption{StereoSGBM parameters used in the reported run}
  \label{tab:sgbm-params}
  \begin{tabular}{lr}
    \toprule
    Parameter & Actual value \\
    \midrule
    \texttt{window\_size} & $3$ \\
    \texttt{minDisparity} & $16$ \\
    \texttt{numDisparities} & $96$ \\
    \texttt{blockSize} & $5$ \\
    $P_1$ & $72$ \\
    $P_2$ & $288$ \\
    \texttt{disp12MaxDiff} & $1$ \\
    \texttt{uniquenessRatio} & $10$ \\
    \texttt{speckleWindowSize} & $100$ \\
    \texttt{speckleRange} & $32$ \\
    \bottomrule
  \end{tabular}
\end{table}

\subsection{Disparity, Valid-Disparity Mask, and Relative Depth}

The OpenCV matcher returns integer disparities with a fixed scale factor. The program divides
these values by $16$ to obtain a floating-point array in pixel units. In this run, a disparity is
considered valid when it is finite and satisfies $d>15$; all other positions are stored as
\texttt{NaN} in \codepath{disparity_raw.npy}, so they do not enter the range, proportion, or
relative-depth statistics. The core logic is:

\begin{lstlisting}[language=Python]
disparity = matcher.compute(left_match, right_match).astype(np.float32) / 16.0
valid = np.isfinite(disparity) & (disparity > 15)
disparity_raw = np.where(valid, disparity, np.nan)
depth_proxy = np.where(valid, 1.0 / disparity_raw, np.nan)
\end{lstlisting}

The relative-depth array is computed as $1/d$ and has no physical unit. It is used only to express
relative front-to-back ordering. To generate a stable display image, the program takes the first
and 99th percentiles of the valid pixels as the display range and linearly maps that range to
0--255. Therefore, the grayscale values in the PNG files are not the raw disparities or raw
relative-depth values; quantitative analysis must read the \texttt{npy} arrays instead.

\subsection{Output and Exception Checks}

After writing each file, the program checks that the file exists and is non-empty. The outputs
include \codepath{left.png}, \codepath{right.png}, \codepath{disparity_raw.npy},
\codepath{disparity.png}, \codepath{valid_mask.png}, \codepath{depth_proxy.npy},
\codepath{depth_proxy.png}, and \codepath{params.json}. The parameter summary also records the
input paths, processed dimensions, matcher parameters, Python/OpenCV/NumPy versions, and output
paths, so that the numerical values reported here can be traced back to the same set of result
files.
